% 【基础正文】所属：第3章（由 chapters/revised/03-justification-context.tex \input 引入）
\minitopic{闭合与演绎扩展}知识的\term{闭合原则}{closure principle}通常表述为：若主体知道$p$，知道$p$蕴涵$q$，并据此推出$q$，则主体知道$q$。闭合解释知识为何能经证明扩展，也是怀疑论论证的枢纽。德雷茨基和诺齐克为保存普通知识而限制闭合\parencite{dretske1970,nozick1981}；反对者认为，这会允许主体知道“这是斑马”，却不知道“这不是伪装得完美的骡子”。 闭合不等于“已知命题的一切逻辑后果都已被知道”。主体必须认识蕴涵并完成推理；否则一个人会自动知道其数学信念所有遥远后果。还要区分知识是否闭合与确证是否闭合，两者可能有不同答案。

\minitopic{语境主义}语境主义主张，知识归属句“$S$知道$p$”的真值条件随归属者的会话语境改变\parencite{derose1992,lewis1996}。普通语境只要求排除日常错误；怀疑语境把欺骗情景变得显著，标准因而上升。变化的是“知道”的适用标准，不一定是主体的证据或心理状态。 其优点是同时解释普通断言与怀疑论压力。困难则有三点：我们是否真能听出“知道”像“高”一样具有语境敏感性？为什么怀疑情景一经提及便改变真值？低标准语境中的知识是否被贬得太容易？

\minitopic{主体敏感与实践侵入}银行案例也可得到另一解释：风险改变的是主体本人是否知道，而非归属者使用“知道”的语义。\term{实践侵入}{pragmatic encroachment}认为，实践利害可以影响知识或确证门槛\parencite{stanley2005,hawthorne2004}。在高风险情境，同样概率证据不足以合理行动，因而也不足以构成知识。 反对者担心知识会变得过分不稳定：银行营业时间不会因房贷风险改变。支持者回应，事实不变，但证据是否足以承担错误代价会变。这里应区分三种主张：风险改变应当复核的行动规范；风险改变相信的理性；风险直接改变知识。前两项不自动推出第三项。

\begin{argumentbox}
实践侵入论证的一种形式：若知道$p$，就可以把$p$当作行动中的既定前提；在高风险情境下，主体不可以仅凭当前证据把$p$当作既定前提；所以主体不知道$p$。争点集中在第一前提：知识是否总足以许可行动？
\end{argumentbox}

\begin{comparison}
中国思想中的“权”常涉及在具体情境中衡量行动，而非改变事实真值。它可帮助我们看到，认识判断嵌在实践中；但不能据此直接推出语义语境主义。把行动上的审慎、相信的规范和“知”的真值条件分层，是比较能够产生分析收益的地方。
\end{comparison}

\minitopic{三种变化不能混淆}第一是\textbf{证据变化}：银行客户新得知周六可能因假日停业，认识位置实际变差。第二是\textbf{话语变化}：证据不变，但谈话把错误可能显著化，“知道”的归属标准升高。第三是\textbf{利害变化}：证据与话语内容都可不变，只是判断错误的代价提高。案例描述若同时改变三项，就不能证明究竟哪项影响知识。 设计思想实验时应逐项控制。低风险与高风险主体应具有相同记忆、相同银行公告和相同推理能力；归属者语境也应固定。然后询问知识判断是否仍变化。即使判断变化，还需排除读者只是把“应当再查一次”误报成“不知道”。实验直觉需要理论诊断，而不是直接投票。

\minitopic{断言、行动与知识}知识常被提出为断言规范：只有知道$p$时才应直接断言$p$。它也被提出为行动理由规范：只有知道$p$时才可把$p$作为不再审查的行动前提。两种规范解释了为何高风险会反照到知识判断，但都存在例外压力。紧急情况下，低于知识的高概率信息也许足以行动；教学或推测语境中，明确保留语气的断言也不要求知识。 更精细的做法是区分直接断言“$p$”、证据报告“迹象支持$p$”和置信表达“我很有把握$p$”。数字与科学交流尤其需要这种层级。把所有合理言说都压成知识断言，会丢失不确定性；把知识完全与行动分离，又难解释知识为何具有实践意义。

\minitopic{案例工作坊：医疗复核}同一张清晰影像由医生判断为良性。若只是常规体检，患者接受一年后复查；若手术明天进行且错误会致命，团队要求第二位医生复核。影像证据并未变，复核要求却变强。我们可以说医生原先知道而现在仍知道，只是行动规范提高；也可以说高风险使当前证据不再足以构成知识。两种描述对实际行动给出相同建议，却对知识与实践关系给出不同理论。 判断哪种解释更好，应观察其他语言行为：医生是否愿意说“影像显示良性，但术前仍须确认”？这句话既保留证据结论，也承认行动门槛。现实沟通常通过限定语管理风险，未必需要让“知道”的真值频繁翻转。

\minitopic{风险论证必须标出中间桥梁}从“错误后果严重”不能直接推出“主体不知道”。可能的桥梁至少有三条：高风险使更多错误可能成为相关替代；高风险改变主体应拥有的证据数量；高风险只改变能否以该命题为行动前提。三条分别属于知识条件、确证强度与行动规范，反例也不同。分析者应先选择桥梁，再说明为何后果会通过它影响评价。若无法给出桥梁，最稳妥结论只是需要复核，而非知识已经消失。
